\documentclass[10pt,a4paper]{amsart}
\usepackage[margin=19mm]{geometry}
\usepackage{amsmath,amssymb,mathtools}
\usepackage[scr=rsfs,cal=euler]{mathalfa}
\usepackage{xfrac}
\usepackage[unicode,hidelinks,bookmarks=false]{hyperref}
\newcommand{\ie}{i.e.\ }
\newcommand{\cf}{cf.\ }
\newcommand{\resp}{respectively\ }
\newcommand{\Subsection}{\subsection}
\providecommand{\nobreakdash}{\nobreak\hbox{-}}
\setlength{\emergencystretch}{2em}
\multlinegap0pt
\usepackage{float}
\usepackage{amssymb}
\usepackage{enumerate}
\usepackage{multirow}
\usepackage{mathtools}
\usepackage[shortlabels]{enumitem}
\usepackage{caption}
\newtheorem{theorem}{Theorem}
\newcommand{\LP}{\left(}
\newcommand{\R}{\mathbb R}
\newcommand{\RP}{\right)}
\title{Polynomial parametrization of some interesting knotted surfaces}
\author{Louis H. Kauffman, Tumpa Mahato, Rama Mishra}
\date{Source: arXiv:2507.19448v2 (CC BY 4.0)}
\begin{document}
\maketitle

\noindent\emph{Source and licence.} Polynomial parametrization of some interesting knotted surfaces. Version arXiv:2507.19448v2 (CC BY 4.0), distributed by arXiv under the Creative Commons Attribution 4.0 International licence, http://creativecommons.org/licenses/by/4.0/, as recorded in the arXiv licence metadata for this exact version and shown on the abstract page https://arxiv.org/abs/2507.19448v2. Full licence notice: \texttt{licenses/arxiv-2507.19448-CC-BY-4.0.txt}.\par\smallskip

\noindent\textit{Editorial scope.} Counted unit: the article's theorem theorem (source0.tex lines 681--688). The article's complete original proof of it is delivered with it (source0.tex lines 690--722, one \texttt{proof} environment of 33 lines). No mathematics is written, repaired or re-derived: the statement and the proof are the article's own lines, in the article's own notation, and no retained line is substituted or re-flowed.  No further source-local context is needed: every cross-reference of every retained line resolves inside the two delivered spans.  One declared adaptation: exactly 1 ancillary strings inside the retained spans are omitted (image-inclusion commands and, where the source repeats a label, the repeated label definitions) and no other line differs from the source; the figure environments, with their placement, captions and labels, are kept, and the package records the omitted strings in fidelity receipts bound to the affected bodies.  No retained line contains a citation, so the wrapper carries no bibliography and no bibliography entry.  The retained lines contain the article's own diagram code, which the wrapper typesets with the standard tikz packages; no image file is delivered and the package declares no asset dependency.  Editorial formatting only, in addition: an AMS \texttt{amsart} preamble that restates the article's own declarations and macros used by the retained lines, namely the environments \texttt{theorem} on separate counters, together with the standard packages those lines need.  The retained environments print in this document as Theorem 1 in order of appearance, whereas the article prints the same items with the numbering of its own full document.  The complete native source of the article as distributed in the arXiv e-print for this exact version is bundled unchanged as \texttt{sources/182211/source0.tex}.
\begin{theorem}
	Given a classical knot $K$ with given trigonometric  parametrization $\phi': [0,2\pi] \rightarrow \R^3$ given by
	\[\phi'(t)=\LP f(t),g(t),h(t)\RP ,\]
	there exists functions $S_{c}(t)$ and $S_{w}(t)$ such that 
	the image of $[0,2\pi]\times [0,2\pi]$ under the map $\tilde{\Psi}: \R^2 \to \R^4$ defined by
	\[\tilde{\Psi}(t,s)=\LP f(t),\,(r-d_{1}S_{c}(t))\,g(t)+\cos s,\,(r-d_{1}S_{c}(t))\,g(t)+\sin s+d_{2}S_{w}(t),\, h(t)\RP,\]
	where $r, d_{1}, d_{2} \in \R^{+}$, is isotopic to the $Tube(K)$.
\end{theorem}

\begin{proof}
	Take the projection of $K$ on $XY$ plane given by 
	$\{f(t),g(t) \, \vert \, t \in[0,2\pi]\}$
	and let there be $n$ number of classical crossings, denoted by $c_{1},c_{2},\cdots,c_{n}$, and $m$ number of welded crossings, denoted by $w_{1},w_{2},\cdots,w_{m}$. Additionally, let $\{t_{i},s_{i}\}$  be the pair of values of $t$ where $t_{i}$ corresponds to the overcrossing point and $s_{i}$ corresponds to the undercrossing point of $c_{i}$ respectively for $i=1,\cdots, n$. Also, let $\{\bar{t}_{j},\bar{s}_{j}\}$  be the pair of values of $t$ where $\bar{t}_{j}$ corresponds the first encounter of the crossing $w_{j}$ and $\bar{s}_{j}$ corresponds the second encounter of the crossing $w_{j}$ for $j=1,\cdots, m$. \\
	First, putting a circle with a fixed radius $r \in \R_{+}$ at each point of the projection is given by
	\begin{equation*}
		\left\{\big(f(t),\,r \,g(t)+\cos s,\,r\, g(t)+\sin s \big)\, \middle| \, t,\,s \in[0,2\pi]\right\}.
	\end{equation*}
	
	Now, we construct functions for shrinking the radius of the lower tube and displacing the tubes from each other, denoted by $S_{c}(t)$ and $S_{w}(t)$ respectively.\\
	Take intervals of the same length, say $L$, around each $t_{i}$ and $s_{i}$ such that no two intervals intersect. Now, we define the bump function with centre $c$ and width $w$ (Fig. \ref{Fig:bump}) by
	\[B(x,c,w)=
	\begin{cases}
		\exp\left(-\frac{w^2}{w^2-(x-c)^2}\right) & \text{if}\; | x-c| <w \\
		0 & \text{otherwise.} 
	\end{cases}\]
		\begin{figure}[H]
		\centering 
		
		\caption{The bump function $B(x,c,w)$ with center $c=0$ and width $w= \frac{\pi}{2}$}
		\label{Fig:bump}
	\end{figure}
	The function $S_{c}(t)$ is defined by summing the functions $B(x,c,\frac{L}{2})$ with centres at the undercrossing points $s_{i}$ of all the classical crossings. On the other hand, the function $S_{w}(t)$ is defined by taking an alternating sum of $B(x,c,\frac{L}{2})$ with centres at the welded crossing points $\bar{t}_{j}$ and  $\bar{s}_{j}$ of all welded crossings. Thus,
		\begin{align*}
			S_{c}(x)=& \sum_{i=1}^{n} B(x,s_{i},\frac{L}{2}),\\
			S_{w}(x)=& \sum_{j=1}^{m} \LP B(x,\bar{t}_{j},\frac{L}{2})-B(x,\bar{s}_{j},\frac{L}{2})\RP.
		\end{align*}

	\noindent
	The  parametrization of the ribbon torus $Tube(K)$ is given by $\tilde{\Psi}: [0,2\pi]\times [0,2\pi] \rightarrow \R^4$ which is defined by
	\[\tilde{\Psi}(t,s)=\LP f(t),\,(r-d_{c}S_{c}(t))\,g(t)+\cos(s),\,(r-d_{c}S_{c}(t))\,g(t)+\sin(s)+d_{w}S_{w}(t),\, h(t)\RP,\]
	where $r, d_{c}, d_{w} \in \R_{+} $ are the radius of the tube, shrinking factor and the displacement factor, respectively, and $h(t)$ provides the over/under information of the tubes in $\R^{4}$. This completes the proof. 
\end{proof}

\end{document}
